\chapter{Turning Chemical Metadata Into Research Variables}

\WeekHeader{6}{Turning Chemical Metadata Into Research Variables}{Convert enrichment outputs into clean variables for grouping, filtering, and comparison.}{Analysis-ready trait table and project-specific grouping dictionary.}

\section{Why This Matters}
Full sentences are difficult to analyze. Researchers need discrete values: hazard codes, pathway IDs, organism names, property bins, pharmacologic actions, measurement units, reaction roles, and presence or absence of curated traits. Week 6 teaches students to build those variables without inventing unsupported labels.

\section{Concepts}
\begin{itemize}
  \item \textbf{Trait type}: a broad evidence domain such as safety, sensory, pathway, taxonomy, bioactivity, or physicochemical behavior.
  \item \textbf{Trait value}: the specific clean value, such as a GHS code, KEGG pathway, organism, or functional term.
  \item \textbf{Matrix key}: a stable column name used in a wide analysis matrix.
  \item \textbf{Confidence}: a source-aware estimate of how reliable or direct an extracted trait is.
  \item \textbf{Project grouping dictionary}: the student's reviewed mapping from traits to analysis groups.
\end{itemize}

\section{Guided Lab}
Start with the normalized tables:

\begin{rcode}
traits <- result$ChemicalTraits
ontology <- result$ChemicalTraitOntology
matrix <- result$ChemicalTraitMatrix

table(traits$TraitType, useNA = "ifany")
table(ontology$OntologyDomain, useNA = "ifany")

head(traits[order(traits$Query, traits$TraitType), ])
\end{rcode}

Build a focused matrix:

\begin{rcode}
core_matrix <- chemicalTraitMatrix(
  result,
  profile = "core",
  mode = "binary",
  min_confidence = "medium",
  max_traits = 100
)

confidence_matrix <- chemicalTraitMatrix(
  result,
  profile = "full",
  mode = "confidence",
  min_confidence = "medium",
  max_traits = 200
)
\end{rcode}

\section{Grouping Rules}
Every grouping variable must pass four tests:
\begin{enumerate}
  \item It is useful for the research question.
  \item It is extracted from a traceable source field.
  \item It has a reproducible rule.
  \item It does not overstate what the database evidence proves.
\end{enumerate}

\begin{goldbox}{Good Versus Weak Grouping Variables}
\textbf{Good}: A binary variable marking chemicals with source-backed GHS hazard codes. The evidence row stores the code and source.

\textbf{Weak}: A manually assigned variable called ``dangerous'' based on reading broad safety sentences.
\end{goldbox}

\section{Independent Extension}
Students design a project-specific grouping dictionary with 8 to 20 variables. Each row should include:
\begin{itemize}
  \item Variable name.
  \item Source table.
  \item Source field.
  \item Extraction rule.
  \item Allowed values.
  \item Confidence rule.
  \item Reason the variable matters.
\end{itemize}

\section{Deliverables}
\begin{tabularx}{\textwidth}{p{0.25\textwidth}YY}
\DeliverableTableHeader
Trait table & Cleaned long table with selected traits and confidence fields. & CSV\\
Trait matrix & Wide matrix ready for visualization or modeling. & CSV\\
Grouping dictionary & Documents how each project variable was derived. & CSV or spreadsheet\\
\end{tabularx}

\section{Quality Checklist}
\begin{checklistbox}{Before Week 7}
\begin{itemize}
  \item No grouping variable is based only on unsupported intuition.
  \item The student can trace each selected variable to a source table and evidence row.
  \item Matrix columns are named clearly enough to interpret later.
  \item Missing values are handled deliberately.
\end{itemize}
\end{checklistbox}
